\chapter{The Realised Trajectory: Cycles 0 to 5}
\label{ch:trajectory}

\Cref{ch:benchmarks} defined the instruments. This chapter reads them across the run
the system actually performed: six cycles, numbered zero through five, on the
five-benchmark panel. Everything here is measured, not designed. Where the earlier
chapters said what \emph{should} happen, this one says what \emph{did}, and it does
not hide the places where the two diverge.

\begin{intuition}[What a cycle is, and what we are watching]
Cycle zero is the cold start: the full architecture is in place, but the model has
not yet been fine-tuned on its own answers. Each later cycle fine-tunes on the
verified episodes accumulated since the last boundary, checks retention, and re-scores
memory. We watch five things across the six cycles: pooled accuracy and honesty, the
per-benchmark spread, how the four-outcome gate shifts, how work moves between the
three tiers, and what the retention guard does when an update goes wrong. The single
most important habit to keep is the one \Cref{ch:benchmarks} argued for: read accuracy
and honesty on separate axes, because the loop moves them differently.
\end{intuition}

\section{The shape of the run}
\label{sec:traj-shape}

The run committed three cycles, aborted one, and committed a final one. That sentence
is the spine of the whole chapter.

\begin{table}[htbp]
\centering
\small
\begin{tabular}{@{}lccccl@{}}
\toprule
\textbf{Cycle} & \textbf{Pooled EM} & \textbf{Pooled CHM} & \textbf{Worst probe} & \textbf{Floor} & \textbf{Outcome} \\
\midrule
C0 (cold start) & 0.541 & 0.135 & ---   & ---  & baseline \\
C1              & 0.543 & 0.112 & 1.006 & 0.93 & commit \\
C2 (best)       & 0.544 & 0.107 & 0.949 & 0.93 & commit \\
C3              & 0.532 & 0.115 & 0.976 & 0.93 & commit \\
C4              & 0.532 & 0.115 & \textbf{0.886} & 0.93 & \textbf{abort} \\
C5 (final)      & 0.517 & 0.109 & 0.958 & 0.93 & commit \\
\bottomrule
\end{tabular}
\caption[The six-cycle trajectory at a glance]{The realised trajectory. Pooled exact
match and the composite hallucination metric are unweighted means over the five panel
benchmarks (the misconception benchmark scored by the language-model judge). The
\emph{worst probe} column is the lowest of the three retention ratios at that cycle;
the cycle aborts when it falls below the floor $\rhomin = 0.93$, which happened at C4.
C2 is the within-trajectory best cycle, C5 the final one.}
\label{tab:traj-glance}
\end{table}

\begin{figure}[htbp]
\centering
\begin{tikzpicture}
\begin{axis}[
  name=emax, width=12.0cm, height=4.0cm,
  xlabel={}, ylabel={pooled EM}, xmin=-0.3, xmax=5.3, ymin=0.50, ymax=0.60,
  xtick={0,1,2,3,4,5}, xticklabels={C0,C1,C2,C3,C4,C5}, ytick={0.50,0.53,0.56,0.59},
  axis lines=left, font=\small\sffamily, tick label style={font=\scriptsize},
  legend style={at={(0.5,1.28)}, anchor=north, legend columns=2, draw=cbInk!25, font=\scriptsize},
  grid=both, grid style={cbInk!8},
]
  \addplot[cbExample!70!black, line width=1.4pt, mark=*, mark size=1.6pt] coordinates {
    (0,0.571)(1,0.584)(2,0.581)(3,0.570)(4,0.571)(5,0.549)};
  \addlegendentry{three training benchmarks}
  \addplot[cbInk!70, line width=1.4pt, mark=square*, mark size=1.5pt] coordinates {
    (0,0.541)(1,0.543)(2,0.544)(3,0.532)(4,0.532)(5,0.517)};
  \addlegendentry{full five-benchmark panel}
  \draw[cbPitfall!70!black, dashed, line width=1pt] (axis cs:4,0.50) -- (axis cs:4,0.60);
  \node[font=\scriptsize, text=cbPitfall!70!black, anchor=south] at (axis cs:4,0.50) {abort};
\end{axis}
\begin{axis}[
  name=chmax, width=12.0cm, height=4.0cm, at={(emax.south)}, yshift=-1.4cm, anchor=north,
  xlabel={cycle}, ylabel={pooled CHM}, xmin=-0.3, xmax=5.3, ymin=0.10, ymax=0.14,
  xtick={0,1,2,3,4,5}, xticklabels={C0,C1,C2,C3,C4,C5}, ytick={0.10,0.11,0.12,0.13,0.14},
  axis lines=left, font=\small\sffamily, tick label style={font=\scriptsize},
  grid=both, grid style={cbInk!8},
]
  \addplot[cbPitfall!65!black, line width=1.4pt, mark=*, mark size=1.6pt] coordinates {
    (0,0.135)(1,0.112)(2,0.107)(3,0.115)(4,0.115)(5,0.109)};
  \draw[cbPitfall!70!black, dashed, line width=1pt] (axis cs:4,0.10) -- (axis cs:4,0.14);
\end{axis}
\end{tikzpicture}
\caption[The headline trajectory]{The actual deployed trajectory. Top: pooled exact
match holds in a narrow band (full panel $0.517$ to $0.544$), with the training
benchmarks running a few points higher. Bottom: the hallucination metric drops sharply
from the cold start to the first committed cycle and then holds low. The dashed line
marks C4, where the retention guard aborted; its eval ran on the rolled-back C3 model,
which is why C3 and C4 read almost identically.}
\label{fig:traj-headline}
\end{figure}

\begin{bythenumbers}[Read the two axes separately]
On the accuracy axis the run is almost flat: pooled exact match stays between $0.517$
and $0.544$ across all six cycles, peaking at C2. On the honesty axis the run does real
work: the hallucination metric falls from $0.135$ at the cold start to $0.107$ at C2
and holds near $0.11$ thereafter, a reduction of roughly a fifth that the accuracy
trace alone would never reveal. This is the single most important reading of the whole
trajectory: the self-improvement loop is a honesty engine, not an accuracy engine, and
a chapter that reported only exact match would have called it a failure.
\end{bythenumbers}

\section{Accuracy, benchmark by benchmark}
\label{sec:traj-perbench}

The flat pooled accuracy hides a real split, and the split is exactly along the
training-versus-transfer line that \Cref{ch:benchmarks} drew.

\begin{figure}[htbp]
\centering
\begin{tikzpicture}
\begin{axis}[
  width=12.4cm, height=6.0cm,
  xlabel={cycle}, ylabel={exact match}, xmin=-0.3, xmax=5.3, ymin=0.25, ymax=0.82,
  xtick={0,1,2,3,4,5}, xticklabels={C0,C1,C2,C3,C4,C5},
  axis lines=left, font=\small\sffamily, tick label style={font=\scriptsize},
  legend style={at={(1.02,0.5)}, anchor=west, draw=cbInk!25, font=\scriptsize},
  grid=both, grid style={cbInk!8},
]
  \addplot[cbExample!70!black, line width=1.3pt, mark=*, mark size=1.4pt] coordinates {
    (0,0.763)(1,0.740)(2,0.717)(3,0.743)(4,0.743)(5,0.717)};
  \addlegendentry{commonsense (train)}
  \addplot[cbScenario!75!black, line width=1.3pt, mark=triangle*, mark size=1.7pt] coordinates {
    (0,0.610)(1,0.647)(2,0.653)(3,0.640)(4,0.640)(5,0.640)};
  \addlegendentry{multi-step (transfer)}
  \addplot[cbFormula!70!black, line width=1.3pt, mark=square*, mark size=1.4pt] coordinates {
    (0,0.487)(1,0.533)(2,0.533)(3,0.503)(4,0.507)(5,0.477)};
  \addlegendentry{fact-check (train)}
  \addplot[cbIntuition!75!black, line width=1.3pt, mark=diamond*, mark size=1.7pt] coordinates {
    (0,0.463)(1,0.480)(2,0.493)(3,0.463)(4,0.463)(5,0.453)};
  \addlegendentry{factoid (train)}
  \addplot[cbPitfall!70!black, line width=1.3pt, mark=otimes*, mark size=1.6pt] coordinates {
    (0,0.380)(1,0.317)(2,0.323)(3,0.310)(4,0.307)(5,0.300)};
  \addlegendentry{misconception (transfer)}
\end{axis}
\end{tikzpicture}
\caption[Per-benchmark exact match across the run]{The actual per-benchmark exact match
across the six cycles (misconception benchmark scored by the language-model judge).
Four of the five benchmarks hold steady or rise; the misconception benchmark falls
monotonically from $0.380$ to $0.300$, which is the one trace that drags the pooled
average down. \Cref{sec:traj-honesty} shows why that fall is honesty, not failure.}
\label{fig:traj-perbench}
\end{figure}

\begin{precise}[Four steady, one declining]
Against the zero-shot floor of \Cref{ch:benchmarks}, four of the five benchmarks
improve or hold. The open-text factoid benchmark rises from a floor of $0.293$ to a
trajectory band around $0.46$ to $0.49$; the fact-checking benchmark rises from $0.453$
into the $0.48$ to $0.53$ band; the commonsense benchmark holds high, in the $0.72$ to
$0.76$ band, a little above its $0.703$ floor; the implicit multi-step benchmark holds
around $0.64$. The single exception is the misconception benchmark, which falls every
cycle, from $0.380$ to $0.300$, ending below its own $0.433$ zero-shot floor. That one
benchmark is the entire reason the pooled five-benchmark accuracy looks flat while the
three training benchmarks sit a few points higher. The honest move is not to drop the
benchmark but to read its decline with the right metric, which the chapter does in
\Cref{sec:traj-honesty}.
\end{precise}

\section{How the gate shifts}
\label{sec:traj-decisions}

The four-outcome gate of \Cref{ch:composite} is not static. As the model improves, the
mix of store, defer, discard, and abstain moves, and the direction of the movement is
itself a measurement of the system's growing confidence.

\begin{figure}[htbp]
\centering
\begin{tikzpicture}
\begin{axis}[
  width=12.0cm, height=5.4cm, ybar stacked, bar width=15pt,
  symbolic x coords={C0,C1,C2,C3,C4,C5}, xtick=data,
  ymin=0, ymax=1.0, ylabel={share of eval queries},
  font=\small\sffamily, tick label style={font=\scriptsize},
  legend style={at={(0.5,1.18)}, anchor=north, legend columns=4, draw=cbInk!25, font=\scriptsize},
  grid=none,
]
  \addplot[draw=cbExample!70!black, fill=cbExample!30] coordinates {
    (C0,0.25)(C1,0.44)(C2,0.47)(C3,0.56)(C4,0.55)(C5,0.42)};
  \addlegendentry{store}
  \addplot[draw=cbNumbers!75!black, fill=cbNumbers!28] coordinates {
    (C0,0.09)(C1,0.32)(C2,0.34)(C3,0.26)(C4,0.26)(C5,0.39)};
  \addlegendentry{defer}
  \addplot[draw=cbInk!55, fill=cbInk!14] coordinates {
    (C0,0.32)(C1,0.10)(C2,0.07)(C3,0.06)(C4,0.07)(C5,0.02)};
  \addlegendentry{discard}
  \addplot[draw=cbIntuition!75!black, fill=cbIntuition!25] coordinates {
    (C0,0.34)(C1,0.14)(C2,0.12)(C3,0.12)(C4,0.12)(C5,0.17)};
  \addlegendentry{abstain}
\end{axis}
\end{tikzpicture}
\caption[The four-outcome distribution over the run]{The actual decision distribution
across the six cycles, pooled over the five-benchmark fold. At the cold start the gate
is cautious: a third of queries are discarded and a third abstained. By the third cycle
the store and defer shares dominate (store peaks near $0.56$) while discard collapses
from $0.32$ to about $0.06$, the signature of a model the gate increasingly trusts.}
\label{fig:traj-decisions}
\end{figure}

\begin{precise}[From cautious to decisive, and the per-outcome ordering]
At the cold start the gate parks or refuses two thirds of its answers: store $0.25$,
defer $0.09$, discard $0.32$, abstain $0.34$. By the third committed cycle it has
turned decisive: store $0.56$, defer $0.26$, discard $0.06$, abstain $0.12$. Two
movements drive this. The discard share collapses from about a third to a sixteenth as
the model produces fewer answers that are low-confidence yet grounded enough not to
abstain; and the store-plus-defer mass grows as more answers clear the trust
thresholds. The per-outcome accuracy stays strictly ordered at every cycle, stored
answers most accurate and abstentions least, which is exactly the ordering a
calibrated gate must produce. The final cycle relaxes slightly, store $0.42$ and defer
$0.39$, as the post-abort model becomes a little less certain, a movement the next
sections explain.
\end{precise}

\section{The three tiers in motion}
\label{sec:traj-tiers}

The clearest evidence that the loop is doing something real is not in the pooled
accuracy but in how the three tiers perform and how work moves between them.

\begin{figure}[htbp]
\centering
\begin{tikzpicture}
\begin{axis}[
  width=12.0cm, height=5.4cm,
  xlabel={cycle}, ylabel={tier exact match}, xmin=-0.3, xmax=5.3, ymin=0.45, ymax=0.64,
  xtick={0,1,2,3,4,5}, xticklabels={C0,C1,C2,C3,C4,C5}, ytick={0.45,0.50,0.55,0.60},
  axis lines=left, font=\small\sffamily, tick label style={font=\scriptsize},
  legend style={at={(0.97,0.97)}, anchor=north east, draw=cbInk!25, font=\scriptsize},
  grid=both, grid style={cbInk!8},
]
  \addplot[cbIntuition!75!black, line width=1.5pt, mark=*, mark size=1.6pt] coordinates {
    (0,0.602)(1,0.585)(2,0.568)(3,0.566)(4,0.564)(5,0.537)};
  \addlegendentry{Tier 2: zero-shot}
  \addplot[cbScenario!75!black, line width=1.5pt, mark=square*, mark size=1.5pt] coordinates {
    (0,0.482)(1,0.502)(2,0.514)(3,0.503)(4,0.504)(5,0.497)};
  \addlegendentry{Tier 3: retrieval}
\end{axis}
\end{tikzpicture}
\caption[Per-tier exact match: the distillation receipt]{The actual per-tier exact
match across the run. The cheap zero-shot tier (Tier 2) answers its assigned queries
\emph{more} accurately than the expensive retrieval tier (Tier 3) at every single
cycle, opening at a twelve-point gap ($0.602$ against $0.482$) at the cold start. The
gap narrows as the loop runs but never reverses, which is the inference-time signature
of self-improvement folding retrieval-found knowledge into the model's own parameters.
The memory-direct tier (Tier 1) is omitted here: it fires on about one query in a
thousand and returned the correct stored answer every time it did.}
\label{fig:traj-tiers}
\end{figure}

\begin{bythenumbers}[The selection effect, stated honestly]
The per-tier reading must be stated with care, because it is partly a selection effect.
The router sends a query to the zero-shot tier precisely because it judges the query
answerable from parameters, so that tier's queries are easier on average. The honest
claim is therefore the conditional one: \emph{the zero-shot tier answers the queries it
is assigned at least as accurately as the retrieval tier answers the queries it is
assigned}, at every cycle, while costing a fraction as much. It is not the claim that
retrieval is worse in general. Even stated conditionally the receipt is strong: a tier
that needs no retrieval matching a tier that does is exactly what a system that has
distilled its retrieved knowledge into its weights should show. The memory-direct tier
fires on roughly one query in a thousand because the evaluation fold is content-hash
disjoint from everything the system trained or stored on, so near-verbatim repeats are
rare by construction; when it does fire it is a precision-first cache, correct on every
realised firing.
\end{bythenumbers}

\section{Retention, and the cycle that was thrown away}
\label{sec:traj-retention}

Three cycles committed cleanly. The fourth did not, and what the system did about it is
the safety property the whole architecture exists to guarantee.

\begin{precise}[The cycle-four abort]
At the fourth cycle boundary the retention guard ran its three probes against their
cold-start anchors. The general-knowledge probe came back at $1.018$, slightly
improved; the commonsense probe at $0.982$; but the held-out open-text factoid probe
fell to $0.886$, below the floor $\rhomin = 0.93$. Because the guard aborts on the
worst probe, the cycle was thrown away: the freshly fine-tuned adapter was discarded
and the model rolled back to its third-cycle weights. This is why C3 and C4 read almost
identically in \Cref{tab:traj-glance}: the C4 evaluation was run on the restored C3
model. The fifth cycle then re-attempted training from that last good adapter and
committed. The retention figure of \Cref{ch:cycle} draws this firing.
\end{precise}

\begin{intuition}[Why the memory did not roll back with the weights]
The abort reverts the weights but \emph{not} the memory. Every episode admitted during
the fourth cycle survived, and the store kept growing across the boundary. This
asymmetry is deliberate and is the core safety property: memory growth is monotone and
high-precision by construction, since each episode passed the gate, so it is always
safe to keep, whereas a weight update is a global change with no such guarantee, so it
is the thing thrown away. A rollback therefore costs a cycle's training but never a
cycle's remembered knowledge. The reconsideration of deferred episodes is one of the
steps gated on committing, so it is skipped on an abort, which is correct: re-scoring
through unchanged weights would be a no-op.
\end{intuition}

\section{The honesty lens on the misconception benchmark}
\label{sec:traj-honesty}

The one benchmark whose accuracy fell every cycle is also the one that proves the
metric suite of \Cref{ch:benchmarks} was necessary. Read with exact match alone, it
looks like a regression. Read with the honesty axis, it is the opposite.

\begin{bythenumbers}[Falling accuracy, rising honesty]
On the misconception benchmark the judged exact match falls from $0.380$ at the cold
start to $0.300$ at the final cycle. Over the same cycles its abstention rate stays
high, in the band of roughly a quarter to a third of all questions (one hundred of
three hundred at the cold start, sixty-nine at the end), where the zero-shot floor
abstains on none. And its hallucination metric \emph{falls}, from about $0.150$ to
about $0.104$. The three numbers together tell one story: the declining exact match
tracks an increase in honest abstention, not a loss of knowledge. The benchmark is
built so the fluent answer is a common misconception, and a dense retrieval substrate
tends to amplify that framing rather than refute it, so the system increasingly
declines rather than commit to a confident falsehood. A bare model is not licensed to
abstain, pays no abstention penalty, and scores a higher strict exact match by stating
the misconception outright. Exact match rewards that; the honesty axis does not.
\end{bythenumbers}

\begin{defence}[If an examiner asks: ``isn't a falling benchmark just a worse system?'']
Not on this benchmark, and the separation of axes is what lets us say so without
hand-waving. The misconception benchmark is adversarial by design: its questions are
written so the most fluent completion is false. A system that answers all of them
scores a higher exact match precisely by asserting falsehoods confidently, which is the
behaviour the whole project is built to suppress. \caem{} instead abstains on a quarter
to a third of them, so its exact match falls while its confident-wrongness metric on
the same benchmark falls too. The architectural remedy for the lost accuracy is the
retrieval-utility classifier of \Cref{ch:ruc}, which routes these queries away from the
amplifying retrieval tier; the trajectory here is the un-routed behaviour the classifier
was built to correct, reported honestly rather than quietly dropped.
\end{defence}

\section{Inside the loop, and when it stopped}
\label{sec:traj-loop}

Underneath the eval numbers the loop itself left a clean record: a training pool that
grew, a training loss that fell, a memory that accumulated, and a principled reason it
stopped at five rather than running to its configured ten.

\begin{precise}[Pool, loss, memory, and the stop]
The per-cycle training pool grew as memory accumulated: roughly one hundred episodes at
the first cycle, then about fifteen hundred, three thousand, and five thousand by the
fifth. The final-epoch training loss fell across the committed cycles, from about
$0.26$ to about $0.11$, with the aborted fourth cycle reaching the lowest loss of the
run, about $0.07$, the textbook signature of an over-fit update that the retention guard
correctly caught (\Cref{ch:cycle} draws this). The episodic store grew monotonically
across every boundary, including the aborted one, from a few hundred episodes to roughly
ten thousand by the final cycle, and the per-cycle re-verification pruned a declining
tail, about thirteen percent at the first committed cycle down to about four percent at
the last, as the pool's composition shifted toward higher-confidence survivors. The loop
stopped at the fifth cycle, not the configured tenth, because two of its three
early-stop signals fired together: the storage rate had saturated and general-knowledge
retention was holding, which is enough to halt at the first eligible cycle after the
burn-in.
\end{precise}

\begin{intuition}[Why stopping early is a feature, not a shortfall]
A self-improvement loop that ran forever would eventually overfit its own outputs, and
the fourth cycle is the warning shot: lowest training loss, failed retention. The
early-stop gate reads the same evidence and halts while the system is still healthy,
rather than grinding through five more cycles of diminishing returns and rising
forgetting risk. The realised five cycles are not a truncated ten-cycle plan; they are
the loop recognising that it had extracted the gains its own answers could supply and
stopping before it started to harm itself.
\end{intuition}

\section{What the theory predicted, and what the run did}
\label{sec:traj-theory}

\Cref{ch:theory} proved properties about the memory pool; the trajectory is where those
properties either held or did not. They held, with one honest exception.

\begin{bythenumbers}[The theory met the run]
The data-purity floor held: the calibration-fold pool purity stayed in the band $0.72$
to $0.75$ across every committed cycle, clearing the registered floor near $0.64$ by
eight to eleven points. The self-correction prediction held: the per-cycle prune rate
declined monotonically, from about thirteen percent to about four percent, as wrong
episodes were caught and removed faster than they accumulated. The Tier-1 floor held:
the memory-direct tier returned the correct answer on every one of its seven firings
across nine thousand evaluation queries. The convergence prediction held in the weak
sense the theory allowed: four of five committed transitions kept retention above the
floor and the fifth was caught and reverted. The one property the run only partially
supported is monotone improvement, because the composite is re-fitted each cycle and so
violates the fixed-function premise the proof assumes; the chapter reports this as a
partial receipt rather than a clean confirmation, exactly as \Cref{ch:theory} flagged.
\end{bythenumbers}

The trajectory is the architecture's behaviour with every piece switched on. To price
each piece, \Cref{ch:ablations} switches them off one at a time and measures what the
panel and the metrics of \Cref{ch:benchmarks} report when a component is removed.
